% Replace the current Results/Discussion experimental text with this section after reviewing journal style.
\section{Experimental Evaluation}
\subsection{Experimental Design}
The expanded evaluation retained the Breast Cancer Wisconsin (Diagnostic) dataset, the 80/20 train--test split, the 30--32--2 multilayer perceptron, full-batch SGD with learning rate 0.03, two local epochs, ten synthetic clients, and class-wise Dirichlet partitioning used in the proof-of-concept. The evaluation used five random seeds (0--4) and 50 global rounds for the reported experiments. Four algorithmic conditions were evaluated: centralized training, local-only training, FedAvg, and FedProx. The framework profile was not treated as a fifth competing algorithm because, under ideal connectivity, its optimization component is FedProx; instead, the framework contribution was evaluated through separate dropout, connectivity, communication, and update-noise stress experiments.

Statistical heterogeneity was varied using Dirichlet $\alpha\in\{100,10,1,0.1\}$. Client dropout was evaluated at 0\%, 10\%, 20\%, 30\%, and 40\% under $\alpha=0.1$. Connectivity was represented by controlled synthetic scenarios (ideal, moderate, poor, intermittent, and severe) with progressively lower availability, higher latency, and greater stale-update probability. These scenarios are stress-test abstractions rather than measurements of rural Indian networks. An additional update-noise experiment varied Gaussian update perturbation from 0 to 0.02; this experiment is a utility stress test and does not constitute a formal differential-privacy guarantee.

\subsection{Non-IID Baseline Results}
Under near-IID partitioning ($\alpha=100$), centralized training, FedAvg, and FedProx achieved a mean test Macro-F1 of 0.951 $\pm$ 0.012 across five seeds. Under $\alpha=10$, the corresponding mean was also 0.951 $\pm$ 0.012. At $\alpha=1$, the mean Macro-F1 of FedAvg/FedProx decreased to 0.947 $\pm$ 0.011, while the mean worst-client F1 decreased to 0.629 $\pm$ 0.166. Under the strongest heterogeneity condition ($\alpha=0.1$), FedAvg/FedProx achieved 0.947 $\pm$ 0.011 global Macro-F1 and 0.453 $\pm$ 0.020 worst-client F1. Thus, the principal degradation associated with increasing statistical heterogeneity was observed in client-level performance rather than in the global test metric.

FedAvg and FedProx produced numerically indistinguishable global results in this configuration. This finding should not be interpreted as evidence that FedProx is ineffective generally; rather, the present dataset, full-batch optimization, two local epochs, and selected proximal coefficient ($\mu=0.01$) did not produce a measurable separation between the two methods. The experiment therefore supports the use of FedProx as an established baseline, not a claim of superiority by the proposed framework.

Local-only training produced substantially lower mean client performance. At $\alpha=0.1$, its mean client F1 was 0.515 $\pm$ 0.093, compared with 0.819 $\pm$ 0.112 for the federated methods. This indicates that, in this controlled simulation, collaborative aggregation improved average client-level performance relative to isolated local training.

\subsection{Client-Dropout Results}
At $\alpha=0.1$, increasing the synthetic client-dropout rate from 0\% to 40\% did not produce a monotonic decline in global Macro-F1 in the present experiment. FedAvg achieved 0.947 $\pm$ 0.011 at 0\% dropout and 0.951 $\pm$ 0.010 at 40\% dropout; FedProx showed the same values because the optimization trajectories were numerically identical in this configuration. Worst-client F1 remained approximately 0.453--0.456 across the dropout levels. Communication volume decreased as dropout increased, from approximately 2.88 MB at 0\% dropout to 1.77 MB at 40\% dropout.

These results should therefore be interpreted as evidence that the simulated aggregation process can continue operating when fewer clients participate, rather than as evidence that dropout has no effect in real deployments. The non-monotonic performance pattern is partly attributable to stochastic client selection and should not be generalized beyond the controlled simulation.

\subsection{Connectivity Stress Results}
The connectivity experiment used five synthetic delivery conditions. Under the ideal condition, FedProx achieved 0.947 $\pm$ 0.011 Macro-F1 with approximately 2.88 MB of total communication and 1,000 ms of modeled cumulative latency. Under the severe condition, Macro-F1 was 0.941 $\pm$ 0.012, while modeled cumulative latency increased to approximately 179,100 ms and total communication decreased to approximately 1.42 MB. The severe condition also produced a mean of 69.4 stale updates across the 50 rounds.

The results demonstrate that the simulated connectivity layer can expose explicit communication and latency trade-offs. However, because the network parameters are synthetic and the stale-update mechanism is an abstract delivery model, these measurements do not establish the behavior of the framework over actual rural Indian network traces.

\subsection{Update-Noise Utility Stress Test}
The Gaussian update-noise experiment did not show a monotonic reduction in global Macro-F1 over the tested range. At zero noise, Macro-F1 was 0.947 $\pm$ 0.011; at noise standard deviation 0.01 it was 0.949 $\pm$ 0.005; and at 0.02 it was 0.964 $\pm$ 0.008. Worst-client F1 at 0.02 decreased to 0.440 $\pm$ 0.063. The non-monotonic global response is consistent with the small synthetic evaluation and should not be interpreted as a privacy benefit. In particular, the injected Gaussian noise is not accompanied by clipping/accounting sufficient to establish $(\epsilon,\delta)$-differential privacy.

\subsection{Discussion of Experimental Findings}
The expanded evaluation provides three bounded findings. First, global test performance remained relatively stable while worst-client performance deteriorated as statistical heterogeneity increased, supporting the need to evaluate client-level fairness rather than relying exclusively on global metrics. Second, synthetic dropout and connectivity stress changed communication and latency characteristics while leaving global predictive performance within a relatively narrow range in this dataset; this supports the feasibility of evaluating resilience as a systems property but does not constitute field validation. Third, FedAvg and FedProx were not separated by the present hyperparameterization, so no superiority claim is warranted.

The results also clarify the scope of the framework contribution. The experiments validate an integrated simulation architecture combining established federated optimization with controlled heterogeneity, client availability, delivery, and noise stressors. They do not demonstrate clinical effectiveness, legal compliance, formal privacy guarantees, or measured rural-network performance. These remain subjects for future validation using real PHC connectivity traces, real healthcare data where ethically and legally permissible, formal privacy accounting, and deployment-oriented evaluation.
